\section{A smaller equation in twelve positive integers}
\label{sec:positive}

Use all twelve coordinates in (2.3) as positive integers. Define
\[
M=kI^2+T^2,\qquad L=1+2HM,\qquad W=b_0+Hr.
\tag{5.1}
\]
The proposed single equation is
\[
\boxed{\quad P_{12}:=N_kN_t\bigl(1+\mathcal N(L,M,W)\bigr)-1=0.\quad}
\tag{5.2}
\]
Together with the finite definitions in Sections~\ref{sec:interface}
and \ref{sec:norm}, this is a fixed integer polynomial, not an
exponential Diophantine equation.

\begin{theorem}\label{thm:positive}
For the complete compiler tuple of Section~\ref{sec:coefficients},
\(P_{12}=0\) has a solution in twelve positive integers if and only if
\(\KTM\) halts on its blank tape. Its exact total degree is \(1081\);
the supplied circuit uses \(116\) binary arithmetic operations.
\end{theorem}

\subsection{Weak bounds before root recovery}

By Lemma~\ref{lem:unit}, a product zero gives
\(\mathcal N=0\) and \(N_k=N_t=\epsilon\in\{1,-1\}\).
The inherited transport bootstrap is unchanged, since it involves only
the outer expressions. It gives
\[
\begin{gathered}
q\ge16,\quad X\ge16,\quad Y\ge4096,\quad C\ge0,\quad
0<R<q^4-2,\quad R>3q\ge48,\\
E\ge q^4>R+2,\qquad
k=R+\epsilon+hE>E,\qquad R<k,\\
kY<c<k(Y+1)<2kY.
\end{gathered}
\tag{5.3}
\]
No square-root recovery, sign of \(\gamma\), or strict inequality
\(C>0\) is assumed in this step.

\subsection{The inner weight}

Write \(a_0=a+2\) temporarily. Since \(X\ge16\), the weak bounds give
\(a_0<2E\). They also give
\[
\ell=kEY>\frac{a_0c}{4},\qquad
\ell>E^2Y>8a_0^2,
\]
and therefore
\[
A>\frac{a_0^2c^2}{16},\qquad A>64a_0^4.
\tag{5.4}
\]
As \(I=ic^2\) with \(i\ge1\), we have
\[
D<a_0^2I^2,\qquad c<I^2,\qquad b=c+RD<2ka_0^2I^2.
\]
Using \(M^2\ge4kI^2T^2\) and \(k>48\), we obtain two separate estimates:
\[
AM^2>\frac{k}{4}Dv^2>8Dv^2,\qquad
AM^2\ge Ak^2I^4>16b^2>8b^2.
\]
Averaging them, and using \(d=Q-1\ge Q/2\), yields
\[
\boxed{\quad dAM^2>2Q(Dv^2+b^2)=2(e+1).\quad}
\tag{5.5}
\]
This strengthening of the inherited constants allows us to use
\(M=kI^2+T^2\) without an extra factor of two.

For either root \(f'=\pm\sqrt D\), any complex \(y'\) satisfying
\(d(y')^2=e-B_af'\) obeys
\[
\begin{split}
d|y'|^2
&\le e+B_a\sqrt D
 <Q(v\sqrt D+b)^2\\
&\le2Q(Dv^2+b^2)<dAM^2.
\end{split}
\]
Thus every such branch satisfies
\[
|y'|<M\sqrt A.
\tag{5.6}
\]
The complex absolute value matters: the right-hand side of the
quadratic equation for a particular branch can be negative.

\subsection{Relative separation of the two main-root branches}

The decisive simplification is the comparison \(A<G\):
\[
\begin{split}
G-A
&>k^2\bigl((E+Y)^2Y^2-EY(EY+1)\bigr)\\
&=k^2\bigl(EY(2Y^2-1)+Y^4\bigr)>0.
\end{split}
\tag{5.7}
\]
Let \(\Phi(x)\in\Z[x]\) be the product over
\(t=\pm\sqrt A,\ f=\pm\sqrt D\) of
\[
d(x-Mt)^2-e+B_af.
\]
Every complex root of \(\Phi\) has the form \(x=Mt+y'\), so (5.6)
gives the uniform bound
\[
|x|<2M\sqrt A<2M\sqrt G.
\tag{5.8}
\]
At \(\mathcal N=0\), one root \(u=\pm\sqrt G\) gives
\(x=(W-Lu)/H\) with \(\Phi(x)=0\).
If \(G\) is not a square, conjugation in \(\mathbb Q(\sqrt G)\)
also makes \(x'=(W+Lu)/H\) a root. But
\[
|x'-x|=\frac{2L\sqrt G}{H}>4M\sqrt G,
\]
contradicting the sum of the two bounds (5.8).
Consequently \(G\) is a square and \(u\) is an integer.
Moreover \(u>0\): if \(u<0\), then
\[
W=Lu+Hx<-L\sqrt G+2HM\sqrt G=-\sqrt G<0,
\]
whereas \(W=b_0+Hr>0\).
We have recovered an integer main root and a rational \(x\).

The earlier absolute separation used an additional factor
\(\ell+k\) in \(L\). Comparing the inner radius with the distance
between the two main-root branches removes that factor.

\subsection{Recovering the remaining integer roots}

For this rational \(x\), put
\[
\mathcal C=d(x^2+M^2A)-e,\qquad
\mathcal V=\mathcal C^2+DB_a^2-4d^2x^2M^2A.
\tag{5.9}
\]
Equation (5.5) implies \(\mathcal C>e+2>0\).
Also \(x\ne0\), since
\[
\Phi(0)=\bigl[(dM^2A-e)^2-DB_a^2\bigr]^2>0,
\]
using \(B_a\sqrt D\le e+1\).
The coherent inverse formulas inherited from the direct merger are
\[
f=-\frac{\mathcal V}{2\mathcal C B_a},\qquad
t=\frac{\mathcal C+B_af}{2dxM},\qquad
y=x-Mt.
\tag{5.10}
\]
Their denominators are nonzero. Squaring the original factor and
eliminating \(t^2=A\) gives
\(\mathcal V+2\mathcal C B_af=0\); substituting back gives
the displayed expression for \(t\).
Thus (5.10) supplies rational numbers with
\(f^2=D\), \(t^2=A\), and
\[
Q(vf-b)^2-(Q-1)y^2=1.
\tag{5.11}
\]
Rational square roots of integers are integers, so \(f,t\in\Z\).
Since \(Q=(\Delta I)^2\) and \(vf-b\in\Z\),
Lemma~\ref{lem:saturation} implies \(y\in\Z\).

Now \(x=Mt+y\) is integral. Reducing \(W=Lu+Hx\) modulo \(H\),
and using \(L\equiv1\pmod H\), gives
\[
\gamma=\frac{u-b_0}{H}\in\Z.
\]
In fact \(\gamma>0\): from \(c>k>E>X\),
\(\Delta>(a+1)^2\), and \(u>0\), we get
\(u>(a+1)c\), hence \(\gamma H>c-X>0\).
The identity
\[
vf-b=c(Tf-1)-Rf^2
\]
uses \(f^2=D\) and restores all four equations (2.8).
Their permitted signed quotient absorbs a negative \(f\) by changing
\((f,T)\) to \((-f,-T)\).
The inherited core theorem forces \(\epsilon=1\) and decodes a
blank halt of \(\KTM\).

Conversely, take a canonical positive normalized compiler witness.
Choose its positive roots \(u,t,y\) and set
\[
r=\gamma+2Mu+Mt+y>0.
\tag{5.12}
\]
Then \(W=Lu+H(Mt+y)\), so one factor in (3.3) vanishes.
The outer factors both equal one. This proves both directions of
Theorem~\ref{thm:positive}.

\subsection{A three-equation version in 110 operations}

Initially the equivalent system is
\[
N_k=1,\qquad N_t=1,\qquad N_0^2=GN_1^2.
\tag{5.13}
\]
Direct comparison in the last equation saves the subtraction
defining \(\mathcal N\).
For the shorter system replace the old transport coordinate by
\(z_{\rm old}=z+1\) and use
\[
\boxed{
\begin{aligned}
N_k&=1,\\
(\kappa+w)C&=F+z(q-1),\\
N_0^2&=GN_1^2.
\end{aligned}}
\tag{5.14}
\]
Every positive solution of (5.14) gives one of (5.13).
Conversely, after the compiler has been recovered, it gives
\(C>0\) and \(X=wq=2^R\) with \(R>3q\). Hence \(w>q>F\), and
\[
(\kappa+w)C-F=(z_{\rm old}-1)(q-1)>0.
\]
Thus \(z_{\rm old}\ge2\), making \(z=z_{\rm old}-1\) positive.
This is the zero-set shift from package 22 \cite{monic22};
it is used after soundness, not as a premature weak bound.
The product (5.2) retains the old coordinate.

The product ledger is \(74\) multiplications and \(42\) additions
or subtractions. The shifted system uses \(72+38=110\) operations;
its residual degrees are \(7,2,1072\).

\subsection{Exact degree certificate}

Give all twelve supplied variables degree one and the four ports
degree zero. The relevant degrees are
\[
\begin{array}{c|rrrrrrrrr}
 &c&H&\Delta&I&D&Q&R&b&e\\ \hline
\deg&5&6&12&11&34&46&4&38&122
\end{array}
\]
and
\[
\begin{array}{c|rrrrrrr}
 &A,G&M&L& W&n&Z_0&Z_1\\ \hline
\deg&22&23&29&11&80&80&40
\end{array}
\]
The second array has a shared \(A,G\) entry.
In \(C_0\) the unique highest homogeneous part is \(-H^2e\),
of degree \(134\), while \(d(Z_0+n)\) has degree at most \(126\).
It follows successively that \(V_0\) has highest part \(H^4e^2\),
of degree \(268\), and \(N_0\) has highest part \(H^8e^4\),
of degree \(536\). In contrast, \(\deg N_1\le488\).
Therefore the final terms have degrees \(1072\) and at most \(998\);
their leading terms cannot cancel.

More explicitly, with \(q_0=\beta J,\ k_0=\eta+\zeta\),
\[
\mathcal N_{\rm top}
=2^{32}i^{48}k_0^{96}w^{80}s^{176}q_0^{656}(q_0-F)^{16}.
\tag{5.15}
\]
The product's highest part is this times
\[
(-hwsq_0^4)\,[w(q_0-F-\alpha)-zq_0].
\]
The coefficient of
\[
J^{677}hzi^{48}\eta^{96}w^{81}s^{177}
\]
is \(2^{32}\beta^{677}\ne0\).
Thus \(\deg\mathcal N=1072\) and \(\deg P_{12}=1081\)
for every authentic tuple, not just a generic small specialization.
